\section{Introduction}

Document parsing converts visually rich document images into structured, machine-readable representations. Moving beyond plain OCR, the task requires preserving not only textual content but also page organization, reading order, tables, formulas, figures, headers, footers, and other layout-dependent elements. Parsed pages are commonly represented in Markdown, so they can be indexed, retrieved, and used by downstream applications~\cite{omnidocbench,puredocbench,paddleocrvl16,mineru25,glmocr}. We study the page-level image-to-Markdown setting, where a model processes a document image and generates a unified Markdown representation of the page.

Current approaches to this task fall into two families. Pipeline methods decompose a page into layout analysis, region-level content recognition, and page-level merging~\cite{paddleocrvl16,mineru25pro,glmocr}. They remain strong on mainstream document parsing leaderboards such as OmniDocBench v1.6~\cite{omnidocbench,omnidocbenchLeaderboard}, where the top three methods at the time of writing are pipeline-based. This design, however, complicates deployment. Layout parsing and content recognition are usually deployed as separate models, often with different runtime loads. Errors also accumulate across stages: missed table boundaries, imprecise formula crops, or wrong reading-order assignments cannot be fully rectified by a downstream recognizer. End-to-end methods take the opposite route, using a single model to read the document page image and generate the Markdown representation in one pass~\cite{hunyuanocr15,unlimitedocr,jiang2026ovisocr}. The one-pass design simplifies deployment and lets the model condition on page-level context throughout generation. Nevertheless, existing end-to-end methods still lag behind pipeline methods in parsing performance on mainstream benchmarks.

We believe the end-to-end approach is more elegant, and develop \ours{} as a compact document parser by post-training Qwen3.5-0.8B, the smallest model in the Qwen3.5 family~\cite{qwen35}. Our goal is to achieve state-of-the-art document parsing while keeping the deployment footprint small, using a single end-to-end model that surpasses the pipeline methods currently leading public leaderboards. Achieving this with a compact backbone requires both a carefully designed data engine and a training recipe suited to long outputs. The data engine combines real documents with synthetic pages. Real-document annotations are produced by specialized parsers, normalized into a unified Markdown schema, and cleaned through rule-based filtering and subset-level spot-checking. Synthetic page generation starts from HTML templates built from challenging samples; the same HTML source is used to render document images and produce Markdown annotations. Training proceeds through supervised fine-tuning, RL training on a 4B branch, on-policy distillation into the 0.8B model, and model fusion.

We evaluate OvisOCR2 on two public benchmarks and an in-house benchmark. On OmniDocBench v1.6~\cite{omnidocbench,omnidocbenchLeaderboard}, OvisOCR2 reaches an overall score of 96.58, setting a new state of the art and moving an end-to-end parser ahead of the leading pipeline methods on this leaderboard. On PureDocBench~\cite{puredocbench}, OvisOCR2 also ranks first with an Avg3 score of 75.06. To complement these public benchmarks, we construct an in-house benchmark of more than 1,000 pages spanning a broader range of document scenarios. OvisOCR2 achieves the highest overall score on this benchmark and leads across all three difficulty tiers. In particular, for handwriting and complex-table documents, which are both challenging to parse yet common in real-world workflows, OvisOCR2 delivers the strongest overall performance on the corresponding subsets.
